\scanpage{29}
\mspage{2}
\margin{28}
\[
\deleted{\mu_{\widetilde Q_{\unsure{a^\vee}}}^{\omega}\mu_{\widetilde Q_a}^{\omega}
=\unsure{\lambda}_{\widetilde Q_a}^{\omega}
=\widetilde c_a^\omega(\unsure{\lambda}_{\widetilde P_a^\omega})}
\qquad
\deleted{\lambda_{\widetilde Q_a}^{\omega}=}
\]
\[
\deleted{\boxed{
\underbrace{\widetilde A_{\omega^{\unsure{m}}r^*}\longrightarrow
\widetilde A_{\omega r^*}\widetilde A_{r^*}}_{\text{déf }\mu_{\widetilde Q_a}^{\omega}}
\widetilde c_r=\widetilde c_r\widetilde d_{\omega^{\unsure{n}}r^*}\cdots
\widetilde d_{\omega r^*}\widetilde d_{r^*}}}
\]
\deleted{conséquence de (1)}

$\mu_{\widetilde Q_a}^{\omega}$ flèche fibre
\uuline{indép. de $\omega$}
\margin{\deleted{de l'autre}\quad $\vee$}

\[
(\unsure{3})\quad\boxed{\mu_{\widetilde R_r}\circ\widetilde b_r
=\widetilde b_{r^\vee}\mu_{\widetilde Q_a}^{\deleted{\unsure{\omega}}}}
\]
\[
(4)\quad\boxed{\mu_{\widetilde R_{r'}}\circ\widetilde a_r
=\widetilde a_{\unsure{r^\vee}\deleted{\unsure{r'}}}\mu_{\widetilde R_r}}
\]
\[
(2)\quad\boxed{\mu_{\widetilde Q_a}\widetilde c_r
=\widetilde c_{r^\vee}\widetilde d_{\omega^{\unsure{n}}r^*}\cdots\widetilde d_{r^*}}
\Longrightarrow\unsure{\mu}
\]
\[
\uline{\text{via (1)}}:\quad
\underline{\mu_{\widetilde Q_a}
=\widetilde A_{\omega^{\unsure{n}}r^*}\cdots\widetilde A_{\omega r^*}\widetilde A_{r^*}}
\]
\[
A_r=\widetilde b_r^{-1}\widetilde a_r^{-1}\widetilde b_{r_1}
\]

\uline{d'où} $\mu_{\widetilde R_r}$ via $(3)$
\[
\underline{\mu_{\widetilde R_r}^{\deleted{\unsure{\omega}}}
=\widetilde B_{\omega^{\unsure{n}}r}\unsure{-|\cdots}\widetilde B_{\omega r}\widetilde B_r}
\]
où
\[
\underline{\widetilde B_r\overset{\text{déf}}{=}\widetilde a_{r''}^{-1}\widetilde b_{r'}\widetilde b_r^{-1}}
\]

et $(4)$ devient une relation en les $\widetilde a_r,\widetilde b_r,\widetilde c_r$, à laquelle il faut ajouter l'indépendance de l'expression de $\mu_{\widetilde Q_a}$ p.r. à $\omega$)

\margin{2 groupes de relations en $\widetilde a_r,\widetilde b_r,\widetilde c_r$ (où \deleted{\unsure{que}} les $\mu$ sont éliminés) \unsure{dépendant d'une relation} $a\in A$ $(\mu_{\widetilde Q_a}^{\omega}=\mu_{\widetilde Q_a}^{\omega'})$ ou $r\in R$ (savoir la relation $(4)$).}

\margin{Mais on voit que ces deux relations sont équivalentes de façon $\pm$ triviale}

On peut considérer $(5)$ comme \uline{définition} des $\lambda_{\widetilde P_a^\omega}$, $(6),(7)$ comme \uline{définition} des $\lambda_{\widetilde Q_a}$ et $\lambda_{\widetilde R_r}$ (en termes de $\widetilde a_r,\widetilde b_r,\widetilde c_r$) donc $(8,\unsure{9}\deleted{\unsure{10}})$ donnent des équations en $\widetilde a_r\ \widetilde b_r\ \widetilde c_r$
